% Generated draft: requires mathematical and source review.
\documentclass[11pt]{article}
\usepackage[margin=25mm]{geometry}
\usepackage{fontspec}
\setmainfont{Latin Modern Roman}
\usepackage{amsmath,amssymb,mathtools,unicode-math}
\setmathfont{Latin Modern Math}
\usepackage{enumitem,needspace,xurl,hyperref}
\setlength{\parindent}{0pt}
\setlength{\parskip}{5pt}
\setlength{\emergencystretch}{3em}
\newcommand{\N}{\mathbb{N}}
\newcommand{\Z}{\mathbb{Z}}
\newcommand{\Q}{\mathbb{Q}}
\newcommand{\R}{\mathbb{R}}
\newcommand{\C}{\mathbb{C}}
\newcommand{\F}{\mathbb{F}}
\newcommand{\A}{\mathbb{A}}
\newcommand{\PP}{\mathbb{P}}
\newcommand{\E}{\mathbb{E}}
\newcommand{\eps}{\varepsilon}
\newcommand{\dd}{\,\mathrm{d}}
\begin{document}
\section*{UnsolvedMath ID: 30000601}
\textbf{Source problem code:} catalogue entry
\subsection{Definitions and mathematical statement}
Let $(G_n)$ be the intended sequence of finite graphs with $|V(G_n)|\to\infty$. In the standard bootstrap model, vertices are initially occupied independently with probability $p$. Occupation is permanent, and an unoccupied vertex becomes occupied when the number of its occupied neighbors reaches the prescribed update threshold. The update rule or neighbor threshold is not specified in the extract. Write $\theta_n(p)=\mathbb P(G_n,p)$ for the intended percolation probability, conventionally the probability that eventually every vertex is occupied.

For $0<\varepsilon<1$, preserve the source's strict-inequality quantile convention:
\[
p_\varepsilon(n)=\inf\{p\in[0,1]:\theta_n(p)>\varepsilon\}.
\]
The two requested threshold-sharpness properties, in the usual nontrivial range $0<\varepsilon<1/2$, are
\[
p_{1-\varepsilon}(n)-p_\varepsilon(n)\longrightarrow0
\quad\text{and}\quad
\frac{p_{1-\varepsilon}(n)}{p_\varepsilon(n)}\longrightarrow1
\qquad(n\to\infty),
\]
with the ratio requiring a positive denominator. The relative-width assertion is stronger than the absolute-width assertion: when the thresholds approach zero, vanishing absolute width alone does not imply a ratio tending to one. The phrase ``in particular'' in the source therefore does not establish an implication from the first displayed limit to the second.

The source literally says ``for every $\varepsilon>0$'', but quantiles at $1-\varepsilon$ and $\varepsilon$ require appropriate probability levels; values $\varepsilon\geq1$ do not give the stated nontrivial finite quantiles. The conventional sharp-threshold range is made explicit rather than treating that wording literally for all positive real numbers. The graph families, bootstrap rule and percolation event must be specified before a universal assertion is determined; no positive theorem for arbitrary growing finite graphs is inferred.
\subsection{Short English statement}
Ask whether the bootstrap-percolation transition window shrinks absolutely, and more strongly relative to its threshold location. The graph family and update rule are missing, and the source's unrestricted epsilon wording and claimed implication need qualification.
\subsection{Sources}
\begin{itemize}
\item Supplied report: \emph{Combinatorics, Probability and Computing} (2007), Oberwolfach Reports, DOI 10.4171/owr/2006/48, pp. 2875--2952; proposed year 2006. \url{https://doi.org/10.4171/owr/2006/48}.
\item Local UnsolvedMath registry, permanent ID 30000601, catalogue entry; created 25 August 2026, updated 16 September 2026. It labels the statement corrected\_verified after comparison with the official report. Its supplied literature triage of 21 August 2026 labels the broad question open, reporting the source's ``wide open in most cases'' qualification and later results for important special bootstrap regimes, rather than a general theorem. No further external reference is supplied; no independent status assessment is made here.
\end{itemize}
\end{document}
